\documentclass[11pt]{article}
\usepackage[utf8]{inputenc}
\usepackage[T1]{fontenc}
\usepackage[margin=1in]{geometry}
\usepackage{amsmath,amssymb,amsthm}
\usepackage{booktabs}
\usepackage{array}
\usepackage[hidelinks]{hyperref}
% generated by gen_numbers.py -- do not edit
\newcommand{\gitcommit}{0168f23713a5}
\newcommand{\seedR}{2}
\newcommand{\seedH}{\tfrac{3}{2}}
\newcommand{\seedT}{\tfrac{1}{2}}
\newcommand{\seedPhi}{3}
\newcommand{\nOneCount}{4096}
\newcommand{\nOneMaxR}{1}
\newcommand{\nOneMaxPhi}{1}
\newcommand{\prodRulePairs}{8}
\newcommand{\prodRuleGeneralPairs}{6}
\newcommand{\prodRuleAllPass}{yes}
\newcommand{\rMaxPowers}{5}
\newcommand{\nMaxPowers}{10}
\newcommand{\lastPowerR}{\tfrac{47}{16}}
\newcommand{\nightRatioNeight}{\tfrac{23}{8}}
\newcommand{\nightIndependent}{\tfrac{23}{8}}
\newcommand{\seedMatA}{\begin{pmatrix} 1 & 1 & 1 & 1 \\ -1 & -1 & 1 & 1 \\ -1 & 1 & 1 & -1 \\ 1 & -1 & 1 & -1 \end{pmatrix}}
\newcommand{\seedMatB}{\begin{pmatrix} -1 & 1 & 1 & -1 \\ 1 & -1 & 1 & -1 \\ -1 & -1 & 1 & 1 \\ -1 & -1 & -1 & -1 \end{pmatrix}}
\newcommand{\seedMatC}{\begin{pmatrix} -1 & -1 & 1 & -1 \\ -1 & -1 & -1 & 1 \\ 1 & -1 & 1 & 1 \\ -1 & 1 & 1 & 1 \end{pmatrix}}
\newcommand{\powersTableRows}{1 & 2 & $2$ & $\tfrac{3}{2}$ & $\tfrac{1}{2}$ & $2$ & yes \\ 2 & 4 & $\tfrac{5}{2}$ & $\tfrac{9}{4}$ & $\tfrac{1}{4}$ & $\tfrac{5}{2}$ & yes \\ 3 & 6 & $\tfrac{11}{4}$ & $\tfrac{21}{8}$ & $\tfrac{1}{8}$ & $\tfrac{11}{4}$ & yes \\ 4 & 8 & $\tfrac{23}{8}$ & $\tfrac{45}{16}$ & $\tfrac{1}{16}$ & $\tfrac{23}{8}$ & yes \\ 5 & 10 & $\tfrac{47}{16}$ & $\tfrac{93}{32}$ & $\tfrac{1}{32}$ & $\tfrac{47}{16}$ & yes \\}
\newcommand{\searchTableRows}{$R$ & 3 & 200 & random & $2$ & 2.0000 \\ $R$ & 4 & 60 & tensor + random & $\tfrac{5}{2}$ & 2.5000 \\ $\Phi$ & 2 & 400 & random & $3$ & 3.0000 \\ $\Phi$ & 3 & 200 & random & $\tfrac{103}{39}$ & 2.6410 \\ $\Phi$ & 3 & 120 & tensor + random & $\tfrac{53}{21}$ & 2.5238 \\ $\Phi$ & 4 & 60 & tensor + random & $3$ & 3.0000 \\}
\newcommand{\signedRowsText}{$1$ at $N=2$, $\tfrac{1}{2}$ at $N=4$, $\tfrac{1}{4}$ at $N=6$, $\tfrac{1}{8}$ at $N=8$}
\newcommand{\searchBestPhi}{3}
\newcommand{\searchBestR}{\tfrac{5}{2}}
\newcommand{\searchAnyAboveThree}{none}
\newcommand{\rigorSha}{04b401f06bbf7405}
\newcommand{\searchSha}{281e6038f6e4b1f8}
\newcommand{\zenodoDOI}{10.5281/zenodo.23232389}

\emergencystretch=3em

\newtheorem{theorem}{Theorem}[section]
\newtheorem{proposition}[theorem]{Proposition}
\newtheorem{lemma}[theorem]{Lemma}
\newtheorem{corollary}[theorem]{Corollary}
\newtheorem{conjecture}[theorem]{Conjecture}
\theoremstyle{remark}
\newtheorem{remark}[theorem]{Remark}

\newcommand{\R}{\mathbb{R}}
\newcommand{\Z}{\mathbb{Z}}
\newcommand{\Lam}{\Lambda}
\DeclareMathOperator{\tr}{tr}

\title{The sharp constant of the absolute dyadic triangular Hilbert sum is at least~3}
\author{Xavier Callens (SocrateAI Lab)}
\date{Preprint v1, 2026-10-08}
\hypersetup{pdftitle={The sharp constant of the absolute dyadic triangular Hilbert sum is at least 3}, pdfauthor={Xavier Callens}}

\begin{document}
\maketitle

\begin{abstract}
Let $\Lam(F_0,F_1,F_2)=\sum_{k\in\Z}\sum_{I\in A_k}|L_I(F_0,F_1,F_2)|$ be the \emph{sum of absolute local contributions}
of the dyadic triangular Hilbert form (the stronger quantity bounded in Theorem~1.1 of the recent preprint \emph{An $L^3$ bound for the dyadic triangular
Hilbert form}, OpenAI, 2026, which proves $\Lam\le 40\prod_v\|F_v\|_{L^3(\R_+^2)}$ for real inputs). Let $C^\ast$ be the
best constant in that estimate. We prove $C^\ast\ge 3$. The lower bound concerns the absolute sum only: on our witnesses
the signed per-scale form is much smaller (Section~\ref{sec:signed}), so nothing is claimed for it. The proof is elementary: for inputs constant on dyadic squares, $\Lam$ is a finite
sum of integer-valued traces, and an exact \emph{product rule} describes how it transforms under Kronecker (tensor)
products of inputs. Tensor powers of one explicit $4\times4$ sign pattern $g$ with $(H,T)=(\seedH,\seedT)$ then have ratio
exactly $3-2^{1-r}$, so the supremum is at least~3. All values are exact rationals, computed by two independent
implementations and, for the first instance beyond $5/2$, by a third. We also prove that the sharp constant at
resolution~1 is exactly~1 and that each dyadic level contributes at most $\prod_v\|F_v\|_3$. We conjecture $C^\ast=3$,
and state the conjecture in a sharper, tensor-closed form, $H(F)+3|T(F)|\le3\prod_v\|F_v\|_3$, which is saturated on the
whole tensor family. Exhaustive search at resolution~1 and exact hill-climbs at resolutions 2--4 found no violation.
The gap $3\le C^\ast\le 40$ remains. Nothing here is machine-checked in a proof assistant.
\end{abstract}

\section{Introduction}

The triangular Hilbert transform couples three functions of two variables, each pair sharing one variable, and its
$L^3\times L^3\times L^3$ estimate is a well-known open problem of multilinear harmonic analysis
\cite{thieleproblems,ktzk2015}. A dyadic model, in which the intervals are constrained by binary addition without
carries, was formulated alongside it. The preprint \cite{oai2026dyadic} (one of the manuscripts released in the
\texttt{openai/math} repository, produced by an internal OpenAI model) proves the uniform estimate for this dyadic
model with unrestricted real inputs and the explicit constant~40. That manuscript is not formalised in Lean (no Comparator challenge in that repository states it) and we have
not verified its proof; we use only its \emph{definition} of the form, which we restate in full. Upstream's
\emph{signed} form carries one coefficient $\epsilon_k$ with $|\epsilon_k|\le1$ per scale and is dominated by the
absolute sum; the theorem below is about the absolute sum, the stronger quantity.

The question we address is the size of the best constant. Our contribution is a lower bound.

\begin{theorem}\label{thm:main}
Let $C^\ast$ be the best constant in $\Lam(F_0,F_1,F_2)\le C^\ast\prod_{v}\|F_v\|_{L^3(\R_+^2)}$ over real, bounded,
compactly supported $F_v$. Then $C^\ast\ge3$.
\end{theorem}

The mechanism is a product rule (Theorem~\ref{thm:product}): on inputs constant on dyadic squares, the form of a
Kronecker product $f\otimes g$ is an affine function of the form of $f$ with slope $|T(g)|$, the top term of~$g$. A seed
with $H(g)/(1-|T(g)|)=3$ therefore generates a sequence with ratios $3-2^{1-r}$. The value $5/2$ found earlier by
numerical ascent (Section~\ref{sec:history}) is the second member of that sequence.

\section{Definitions and the finite reduction}

We write $\R_+=[0,\infty)$. For $k\in\Z$ let $D_k=\{[2^k n,2^k(n+1)):n\in\Z_{\ge0}\}$, and for $I\in D_k$ with index
$n_I$ let $h_I=\mathbf 1_{I_{\mathrm{left}}}-\mathbf 1_{I_{\mathrm{right}}}$. Let $\oplus$ be bitwise XOR and
\[
A_k=\{(I_0,I_1,I_2)\in D_k^3:\ n_{I_0}\oplus n_{I_1}\oplus n_{I_2}=0\}.
\]
For real measurable $F_0,F_1,F_2$ on $\R_+^2$ and $I\in A_k$,
\begin{equation}\label{eq:LI}
L_I(F_0,F_1,F_2)=2^{-k}\int_{I_0\times I_1\times I_2}F_0(x_0,x_1)F_1(x_1,x_2)F_2(x_2,x_0)\,h_{I_0}(x_0)h_{I_1}(x_1)h_{I_2}(x_2)\,dx_0\,dx_1\,dx_2,
\end{equation}
and $\Lam(F)=\sum_{k\in\Z}\sum_{I\in A_k}|L_I(F)|$. Theorem~1.1 of \cite{oai2026dyadic} states
$\Lam(F)\le40\prod_v\|F_v\|_3$ for real bounded compactly supported inputs, and extends to $L^3$ by density. We write
$R(F)=\Lam(F)/\prod_v\|F_v\|_3$ and $C^\ast=\sup R$.

\paragraph{Step inputs.}
Fix $N\ge1$, $n=2^N$, and let each $F_v$ be constant on the squares of side $2^{-N}$ inside $[0,1)^2$ and zero
elsewhere; identify $F_v$ with the $n\times n$ matrix $F_v[i,j]$ of its values. Write $T(F)=\int F_0F_1F_2$ over the unit
cube, so that $T=n^{-3}\sum_{i,j,k}F_0[i,j]F_1[j,k]F_2[k,i]=n^{-3}\tr(F_0F_1F_2)$.

\begin{proposition}[finite reduction]\label{prop:finite}
For step inputs at resolution $N$,
\[
\Lam(F)=|T(F)|+H(F),\qquad H(F)=\sum_{l=0}^{N-1}\sum_{I\in A_{-l}}|L_I(F)|,
\]
and every $L_I$ is $2^{l}n^{-3}$ times an integer combination of products of entries. Dyadic dilations leave $R$
unchanged, so restricting the support to $[0,1)^2$ loses nothing.
\end{proposition}

\begin{proof}
At scale $2^k\ge2$ only the triple of intervals containing $[0,1)$ meets the support, all three Haar functions equal $+1$
there, and $L=2^{-k}T$; these scales sum to $|T|\sum_{k\ge1}2^{-k}=|T|$. At scale $2^k\le2^{-N}$ each $h_I$ integrates to
zero on every cell, so $L_I=0$. At level $l$, i.e.\ $k=-l$ with $0\le l\le N-1$, there are $m=2^l$ intervals of $b=n/m$
cells each; writing $s\in\{\pm1\}^b$ for the Haar signs and $F_v^{pq}$ for the $b\times b$ block of $F_v$ on the
interval pair $(p,q)$, the triple $(p,q,p\oplus q)$ contributes
\begin{equation}\label{eq:trace}
L_I=2^{l}n^{-3}\tr\big(SF_0^{pq}SF_1^{q,p\oplus q}SF_2^{p\oplus q,p}\big),\qquad S=\operatorname{diag}(s).
\end{equation}
Dilation invariance follows from the scaling of \eqref{eq:LI}.
\end{proof}

For sign inputs ($F_v[i,j]\in\{\pm1\}$) every $\|F_v\|_3=1$, so $R=|T|+H$ is an exact rational. All computations below
use integer arithmetic on \eqref{eq:trace}; nothing is floating point.

\section{The product rule}

For $f=(f_0,f_1,f_2)$ at resolution $N_1$ and $g$ at resolution $N_2$, let $f\otimes g$ be the input at resolution
$N_1+N_2$ with $(f_v\otimes g_v)[i,j]=f_v[\lfloor i/2^{N_2}\rfloor,\lfloor j/2^{N_2}\rfloor]\,g_v[i\bmod 2^{N_2},j\bmod 2^{N_2}]$
(the Kronecker product: $f$ on the coarse cells, a copy of $g$ inside each). Put
\[
A'(f)=2^{-2N_1}\sum_{p\oplus q\oplus r=0}|f_0[p,q]\,f_1[q,r]\,f_2[r,p]|,
\]
the triangular form of $|f|$ over the admissible coarse triples; $A'(f)=1$ for sign inputs.

\begin{theorem}[product rule]\label{thm:product}
For real step inputs $f,g$,
\[
\Lam(f\otimes g)=|T(f)|\,|T(g)|+|T(g)|\,H(f)+A'(f)\,H(g),\qquad
\|f_v\otimes g_v\|_3=\|f_v\|_3\|g_v\|_3,\qquad T(f\otimes g)=T(f)T(g).
\]
In particular, for sign inputs,
\begin{equation}\label{eq:rule}
R(f\otimes g)=H(g)+|T(g)|\,R(f).
\end{equation}
\end{theorem}

\begin{proof}
Let $F=f\otimes g$, $n_1=2^{N_1}$, $n_2=2^{N_2}$.

\emph{Coarse levels $l<N_1$.} Each Haar function is constant on every coarse cell of side $2^{-N_1}$, so the integral
in \eqref{eq:LI} splits into a sum over coarse cell triples $(a_0,a_1,a_2)$ of
$\pm f_0[a_0,a_1]f_1[a_1,a_2]f_2[a_2,a_0]$ times the integral of $g_0g_1g_2$ over the cell cube, which is
$n_1^{-3}T(g)$ after rescaling the cell to the unit cube. Hence $L_I(F)=T(g)\,L_I(f)$ for every coarse triple, and the
same computation with all signs $+1$ gives $T(F)=T(f)T(g)$.

\emph{Fine levels $l\ge N_1$.} Write $I_v=2^{-N_1}(a_v+J_v)$ with $J_v$ dyadic of level $l-N_1$. The index of $I_v$ is
$a_v2^{l-N_1}+j_v$, so the XOR condition holds iff $a_0\oplus a_1\oplus a_2=0$ and $j_0\oplus j_1\oplus j_2=0$. On the
cube of $I$ the inputs equal $f_0[a_0,a_1]f_1[a_1,a_2]f_2[a_2,a_0]$ times the fine copy of $g$, the prefactor splits as
$2^{l}=2^{N_1}2^{\,l-N_1}$ and the measure contributes $n_1^{-3}$, so
$L_I(F)=f_0[a_0,a_1]f_1[a_1,a_2]f_2[a_2,a_0]\cdot n_1^{-2}\,L_J(g)$. Summing $|L_I(F)|$ over the admissible $(a_0,a_1)$
and over $J$ gives $A'(f)\sum_J|L_J(g)|$ at level $l-N_1$; over all fine levels this is $A'(f)H(g)$.

The norm identity is $\|f_v\otimes g_v\|_3^3=n_1^{-2}n_2^{-2}\sum|f_v|^3\sum|g_v|^3$. Adding the two groups of levels
to $|T(F)|$ gives the formula; for sign inputs all norms and $A'(f)$ equal~1.
\end{proof}

\begin{lemma}[closure and limit]\label{lem:closure}
For sign inputs with $|T|<1$ put $\Phi(g)=H(g)/(1-|T(g)|)$.
\begin{enumerate}
\item[(i)] If $\Phi(f)\le M$ and $\Phi(g)\le M$ then $\Phi(f\otimes g)\le M$.
\item[(ii)] $R(g^{\otimes r})\to\Phi(g)$ as $r\to\infty$, with $R(g^{\otimes r})-\Phi(g)=|T(g)|^{r-1}\big(R(g)-\Phi(g)\big)$.
The sequence is nondecreasing iff $R(g)\le\Phi(g)$, which holds iff $T(g)=0$ or $R(g)\ge1$ (in particular for the seed of
Theorem~\ref{thm:three}); otherwise it is decreasing. Consequently $\sup_rR(g^{\otimes r})=\max(R(g),\Phi(g))$.
\end{enumerate}
\end{lemma}

\begin{proof}
(i) By \eqref{eq:rule} applied twice, $H(f\otimes g)=H(g)+|T(g)|H(f)$ and $|T(f\otimes g)|=|T(f)||T(g)|$. Then
$H(f\otimes g)\le M(1-|T(g)|)+|T(g)|M(1-|T(f)|)=M(1-|T(f)||T(g)|)$.
(ii) By \eqref{eq:rule}, $R_r=H+|T|R_{r-1}$, an affine contraction with fixed point $\Phi$ since $|T|<1$, so
$R_r-\Phi=|T|^{r-1}(R_1-\Phi)$ has a constant sign. Moreover $R_1\le\Phi$ iff $(H+|T|)(1-|T|)\le H$ iff
$|T|(1-H-|T|)\le0$ iff $T=0$ or $R_1\ge1$.
\end{proof}

Thus a tensor power of a single seed gives at most $\max(R(g),\Phi(g))$: beyond the seed's own ratio, what the tensor
construction adds is exactly the limit $\Phi(g)$, and by (i) products of seeds with $\Phi\le M$ never raise $\Phi$ above $M$.

\section{The lower bound}

\begin{theorem}\label{thm:three}
Let $g=(g_0,g_1,g_2)$ be the sign input at resolution~2
\[
g_0=\seedMatA,\quad g_1=\seedMatB,\quad g_2=\seedMatC.
\]
Then $H(g)=\seedH$, $T(g)=\seedT$, $R(g)=\seedR$, and $R(g^{\otimes r})=3-2^{1-r}$ for every $r\ge1$. Hence $C^\ast\ge3$.
\end{theorem}

\begin{proof}
The values of $H$ and $T$ are finite integer computations from \eqref{eq:trace} (Section~\ref{sec:comp}). By
\eqref{eq:rule}, $R_r=\tfrac32+\tfrac12R_{r-1}$ with $R_1=2$, whose solution is $3-2^{1-r}$. Each $g^{\otimes r}$ is a
bounded compactly supported sign input, so $C^\ast\ge\sup_rR(g^{\otimes r})=3$.
\end{proof}

Table~\ref{tab:powers} lists the exact values up to $r=\rMaxPowers$ (resolution $N=\nMaxPowers$), each computed directly
from \eqref{eq:trace} rather than from the recursion, together with the closed form.

\begin{table}[h]\centering
\begin{tabular}{rrllllc}\toprule
$r$ & $N$ & $R(g^{\otimes r})$ & $H$ & $T$ & $3-2^{1-r}$ & match \\\midrule
\powersTableRows
\bottomrule\end{tabular}
\caption{Exact ratios of the tensor powers of the seed $g$ (direct evaluation, integer arithmetic).}\label{tab:powers}
\end{table}

\begin{remark}
The seed $g$ is the sign pattern of the resolution-2 witness found by numerical ascent (Section~\ref{sec:history});
it attains the exhaustive resolution-2 maximum $R=2$ among sign inputs (Section~\ref{sec:comp}). Its tensor square $g\otimes g$ has ratio $5/2$ and explains the value
first reached by floating-point ascent at resolution~4.
\end{remark}

\section{The signed form}\label{sec:signed}

Upstream's form with coefficients is $\Lam_{S,\epsilon}=\sum_k\epsilon_k\sum_{I\in A_k}L_I$. Its supremum over
$|\epsilon_k|\le1$ is $|T|+\sum_{l}\big|\sum_{I\in A_{-l}}L_I\big|\le\Lam$ (scales $k\ge1$ each carry $2^{-k}T$). The product
rule does not transfer: the fine levels of $f\otimes g$ are multiplied by the \emph{signed} coarse sum of $f$, not by
$A'(f)=1$. On the tensor powers of Theorem~\ref{thm:three} this supremum is
\signedRowsText, against absolute ratios $2,\tfrac52,\tfrac{11}4,\tfrac{23}8$. So these inputs give no lower bound for the
signed form, and we state none: the best constant of the signed form lies between $1$ and $C^\ast\le40$ as far as this
paper is concerned.

\section{Upper bounds}\label{sec:upper}

Write $C_N$ for the best constant over step inputs at resolution $N$ (real entries). Clearly $C_1\le C_2\le\dots$ and
$C^\ast=\sup_NC_N$ by density.

\begin{proposition}\label{prop:level}
For real step inputs, $|T(F)|+\sum_{I\in A_0}|L_I(F)|\le\prod_v\|F_v\|_3$ and, for each level $l$,
$\sum_{I\in A_{-l}}|L_I(F)|\le\prod_v\|F_v\|_3$. Consequently $C_N\le N+1$, and $C_1=1$.
\end{proposition}

\begin{proof}
We use the triangular H\"older inequality: for any abelian group $G$ and functions $u_v$ on $G^2$,
$\sum_{x\oplus y\oplus z=0}|u_0(x,y)u_1(y,z)u_2(z,x)|\le\prod_v\|u_v\|_{\ell^3(G^2)}$, because the three maps
$(x,y)\mapsto(x,y),(y,z),(z,x)$ from the admissible set to $G^2$ are bijections, so the generalised H\"older inequality
with exponents $(3,3,3)$ applies and each factor is the full $\ell^3$ norm. (The same holds for the odd class
$x\oplus y\oplus z=1$ and, with Lebesgue measure, for integrals over half-cubes.)

\emph{Level $l$.} Each $|L_I|\le2^{l}\int_{I_0\times I_1\times I_2}|F_0F_1F_2|$. On that cube H\"older with exponents
$(3,3,3)$ gives $\int|F_0F_1F_2|\le\prod_v(2^{-l}\int_{Q_v}|F_v|^3)^{1/3}=2^{-l}\prod_v\|F_v\mathbf 1_{Q_v}\|_3$,
where $Q_0=I_0\times I_1$, $Q_1=I_1\times I_2$, $Q_2=I_2\times I_0$ (the third variable ranges over an interval of
length $2^{-l}$). Summing over $I\in A_{-l}$ and applying the triangular H\"older inequality on the group
$(\Z/2)^l$ of interval indices under XOR, with $u_v(p,q)=\|F_v\mathbf 1_{I_p\times I_q}\|_3$, gives
$\sum_I|L_I|\le\prod_v(\sum_{p,q}\|F_v\mathbf 1_{I_p\times I_q}\|_3^3)^{1/3}=\prod_v\|F_v\|_3$.

\emph{Top terms.} $T\pm L_{(0,0,0)}=\int F_0F_1F_2(1\pm h(x_0)h(x_1)h(x_2))$, and $1+hhh$ equals $2$ on the four
half-cubes of even parity and $0$ otherwise (and $1-hhh$ the reverse). So
$|T|+|L_{(0,0,0)}|=\max(|T+L|,|T-L|)\le2\max_{\text{parity}}\int_{\text{class}}|F_0F_1F_2|$, and on each half-cube
H\"older gives a factor $\tfrac12\prod_v\|F_v\mathbf 1_{Q_v}\|_3$; the four half-cubes of one parity class map
bijectively onto the four quarter-squares of each $F_v$, so the triangular H\"older inequality on $\Z/2$ bounds the sum
by $\tfrac12\prod_v\|F_v\|_3$. Hence $|T|+|L_{(0,0,0)}|\le\prod_v\|F_v\|_3$.

Adding the level bounds gives $C_N\le N+1$; at $N=1$ the only level is $l=0$, so $C_1\le1$, and $F\equiv1$ attains~1.
\end{proof}

The bound $C_N\le N+1$ is attained when all Haar signs are replaced by $+1$ (the Haar-free form), which is how the
numerical instruments are checked for power. The uniform bound $40$ of \cite{oai2026dyadic} exploits cancellation across
levels through a matrix energy $\tr((R^\ast R)^{3/2})$; we have not improved it.

\section{A sharp conjecture}

\begin{conjecture}\label{conj:psi}
For all real step inputs (equivalently all real inputs), $H(F)+3|T(F)|\le3\prod_v\|F_v\|_3$.
\end{conjecture}

Since $R=|T|+H\le H+3|T|$, Conjecture~\ref{conj:psi} implies $C^\ast\le3$, hence $C^\ast=3$ by
Theorem~\ref{thm:three}. The functional $\Psi(F)=H(F)+3|T(F)|$ is the natural object for three reasons. First, for sign
inputs $\Psi\le3$ is equivalent to $\Phi\le3$, so it is exactly the statement that no seed beats the one found. Second,
equality propagates through tensor products: if $\Psi(f)=\Psi(g)=3$ then, by Theorem~\ref{thm:product},
$\Psi(f\otimes g)=H(g)+|T(g)|(H(f)+3|T(f)|)=H(g)+3|T(g)|=3$. The constant input ($H=0,T=1$) and the seed $g$
($H=\tfrac32,T=\tfrac12$) both satisfy $\Psi=3$, so the entire family $g^{\otimes r}$, and every tensor product of
copies of $g$ and $\mathbf 1$, lies on the equality set. Third, at resolution~1 it is a theorem:
$\Psi\le|L|+3|T|\le1+2|T|\le3$ by Proposition~\ref{prop:level}. A proof for general $N$ cannot proceed cell by cell,
because the fine levels only see the admissible (even) coarse triples while $T$ sees all eight, and each square of
$F_v$ is shared by one even and one odd triple; this coupling is where the content of the conjecture lies.

\section{Computations}\label{sec:comp}

All results are in the repository \texttt{xaviercallens/AutoevolveAI} at commit \texttt{\gitcommit}, directory
\texttt{results/openai\_math/hypotheses/tensor/}, produced by \texttt{scripts/openai\_math/hypotheses/tensor/}
(\texttt{h5\_tensor\_rigor.py}, sha256 prefix \texttt{\rigorSha}; \texttt{phi\_search.py}, \texttt{\searchSha}), and
archived at Zenodo (DOI \zenodoDOI).

\paragraph{Independent implementations.} The evaluator in \texttt{h5\_tensor\_rigor.py} was written from
\eqref{eq:LI}--\eqref{eq:trace} and agrees with the earlier evaluator \texttt{h5\_exact.py} on random sign inputs at
resolutions 1--4 (test suite). The value $\nightRatioNeight$ of $w_4\otimes w_4$ (Section~\ref{sec:history}) was obtained
by both, and by a third implementation written by an adversarial verifier agent from the definition, with a different
contraction order ($\nightIndependent$). Controls: the constant input gives exactly~1; the Haar-free form gives exactly
$N+1$; the product rule was checked exactly on \prodRulePairs{} random sign pairs and, in its general form with $A'(f)$,
on \prodRuleGeneralPairs{} random rational pairs (all pass: \prodRuleAllPass).

\paragraph{Exhaustive resolution~1.} Over all $\nOneCount$ sign inputs, $\max R=\nOneMaxR$ and $\max\Phi=\nOneMaxPhi$.

\paragraph{Resolution~2.} The exhaustive maximum of $R$ over the $2^{48}$ sign inputs is $2$. The computation (night lane
of 2026-10-07) fixes $(F_0,F_1)$, reduces by the symmetries $F_v\to-F_v$, and uses that for fixed coefficient vectors
$c_t$ of the six terms, $\max_{F_2\in\{\pm1\}^{16}}\sum_t|\langle c_t,F_2\rangle|=\max_{s\in\{\pm1\}^6}\|\sum_ts_tc_t\|_1$;
each chunk's argmax was rechecked by the exact evaluator. The seed $g$ attains this maximum.

\paragraph{Searches for a seed beating 3.} Exact hill-climbs on $\Phi$ (single and pair sign flips, random and perturbed
tensor starts) and on $R$ found no value of $\Phi$ above 3 (best $\Phi$: $\searchBestPhi$; best $R$: $\searchBestR$;
any above 3: \searchAnyAboveThree). Table~\ref{tab:search} lists the runs. A hill-climb proves nothing; it is reported
as the evidence it is.

\begin{table}[h]\centering
\begin{tabular}{llrlll}\toprule
objective & $N$ & restarts & starts & best (exact) & best \\\midrule
\searchTableRows
\bottomrule\end{tabular}
\caption{Exact sign-flip hill-climbs (\texttt{phi\_search.py}). None exceeds 3.}\label{tab:search}
\end{table}

\section{How the bound was found}\label{sec:history}

The constant was first studied by preregistered numerical experiments in a research pipeline run by language-model
agents. A H\"older-ascent search at resolutions 2--7 reached exactly $2$ at $N\le3$ and $5/2$ at $N\ge4$. The
hypothesis $C^\ast=2$ was preregistered and refuted by its own run; the amended hypothesis $C^\ast=5/2$ survived
fresh-seed runs at $N=6,7$ and was then refuted at $N=8$ by the tensor square of the $N=4$ witness $w_4$, with exact
ratio $\nightRatioNeight$. The product rule was found while explaining that number. Local ascent never found the tensor
extremisers on its own: the recorded ascent runs at resolutions 4--8 ended between $2.49991$ and $2.5$ (the 18 preregistered
ascent chunks of the night lane at $N=5..8$ had extremes $2.4999186$ and $2.49999999995$) and none exceeded $5/2$.
That is consistent with $5/2$ being a local maximum of $R$ among the inputs reachable by small moves, but it does not
show it.

\section{Limitations and AI use}

The proofs are short. Theorem~\ref{thm:product} was read and its identity re-derived numerically by an independent
model verifier; Proposition~\ref{prop:finite}, Lemma~\ref{lem:closure} and Proposition~\ref{prop:level} were written
for this manuscript and have been checked only by the agent that wrote them. None is machine-checked in a proof
assistant. The exhaustive resolution-2 maximum is a computer-assisted claim from a preregistered run; the
resolution-1 maximum was recomputed here. The upper bound $40$ is taken from \cite{oai2026dyadic} and is not verified
here; Theorem~\ref{thm:main} does not depend on it. Whether the true constant is~3 is open.

The computations, the search for the seed, the product rule and this manuscript were produced by Claude agents (Anthropic)
running in Claude Code under a workflow that the author designed and launched: preregistration committed before each
confirmatory run, positive and negative controls, exact arithmetic, and an adversarial model verifier for every lane.
The author chose the problem and the rules. No person has independently refereed this manuscript before this preprint.

\begin{thebibliography}{9}
\bibitem{oai2026dyadic} OpenAI, \emph{An $L^3$ bound for the dyadic triangular Hilbert form}, preprint, 5 October 2026,
in \texttt{github.com/openai/math} (commit \texttt{adc7f124}), \texttt{preprints/An-L3-bound-for-the-dyadic-triangular-Hilbert-form-October-5-2026/}.
\bibitem{ktzk2015} V.~Kova\v{c}, C.~Thiele, P.~Zorin-Kranich, \emph{Dyadic triangular Hilbert transform of two general
functions and one not too general function}, Forum Math.\ Sigma 3 (2015), e25; arXiv:1506.00861.
\bibitem{thieleproblems} L.~Grafakos, D.~Oliveira e Silva, M.~Pramanik, A.~Seeger, B.~Stovall, \emph{Some problems in
harmonic analysis}, arXiv:1701.06637v1 (2017), Section~9, contributed by C.~Thiele.
\end{thebibliography}

\end{document}
